% Source: physical PDF pp. 356--360 (printed pp. 333--337), beginning at
% the Section 31.3 heading on p. 356 and ending immediately before the
% Section 31.4 heading on p. 360.
% Inventory: Eqs. (31.3.1)--(31.3.20), three unnumbered display groups,
% one linked citation marker (reference [11]), and no footnotes, figures,
% tables, or centered three-asterisk dividers.
% Source-visible anomaly preserved: the opening sentence says that the
% weak-field formalism was developed in Sections 31.1--31.3, thereby
% including the present section.
% In accordance with NOTATION.md, the direct product of gauge groups is
% set with a centered dot rather than the multiplication sign in the scan.
% Uncertainties: none.

\subsection{The Gravitino}

In this section we will use the weak field formalism developed in Sections
31.1--31.3 to derive some properties of the gravitino. In particular, by
using an argument of continuity as \(G\to0\) we shall obtain a formula for
the gravitino mass when supersymmetry is spontaneously broken that is valid
to first order in \(G\) but to all orders in all other interactions. (The
original derivation of this formula will be given in Section 31.6.)

First we must verify that the term
\(-\frac12\bar\psi_\mu L^\mu\) in Eq.~\eqref{eq:31.2.9} is the correct
free-field Lagrangian for a massless self-charge-conjugate particle of spin
\(3/2\). The time-honored approach to the construction of suitable
free-field Lagrangians for particles with spin is to guess at the Lagrangian
and then check that it gives a physically satisfactory field equation and
propagator. This has led to some uncertainty for particles of spin \(3/2\)---
for instance, what is usually called the Rarita--Schwinger Lagrangian in
papers on supersymmetry is not the Lagrangian originally proposed by Rarita
and Schwinger.~\hyperref[ch31-ref-11]{[11]} Here we shall follow an approach
in the spirit of Section 6.2: we will first derive the propagator from the
requirements of Lorentz invariance for a massive particle of spin \(3/2\)
and then invert it to find the Lagrangian. We work here with massive
gravitinos for the sake of simplicity and because in the real world we must
take account of the breakdown of supersymmetry, but we will be able to apply
the results obtained in this way to the case of massless gravitinos by
noting that supersymmetry current conservation makes the zero-mass
singularities in the propagator inconsequential.

A spinor field \(\psi^\mu\) with an extra vector index belongs to the
\([(\frac12,0)+(0,\frac12)]\mathbin{\times}(\frac12,\frac12)\)
representation of the homogeneous Lorentz group. To isolate the
\((1,\frac12)+(\frac12,1)\) parts of the free field, we impose the
irreducibility condition
\begin{equation}
  \gamma_\mu\psi^\mu=0\,.
  \tag{31.3.1}\label{eq:31.3.1}
\end{equation}
Rotational invariance and Eq.~\eqref{eq:31.3.1} tell us that the matrix
elements of this field between the vacuum and a massive spin \(3/2\)
particle of momentum \(\boldsymbol q=0\) and spin \(z\)-component \(s\)
will satisfy the conditions
\begin{equation}
  \langle0|\psi^0(0)|s\rangle=0
  \tag{31.3.2}\label{eq:31.3.2}
\end{equation}
and
\begin{equation}
  \sum_{s=-3/2}^{3/2}
  \langle0|\psi^i(0)|s\rangle
  \langle0|\psi^j(0)|s\rangle^*
  \propto\delta_{ij}-\frac13\gamma_i\gamma_j\,,
  \tag{31.3.3}\label{eq:31.3.3}
\end{equation}
with a coefficient that may depend on the rotationally invariant matrix
\(\beta\equiv i\gamma^0\). With the usual Dirac convention that
\(\beta\langle0|\psi^i(0)|s\rangle=\langle0|\psi^i(0)|s\rangle\)
(chosen to simplify the space inversion property of the field) and a
conventional choice of normalization of the \(\beta=+1\) component
analogous to Eq.~\eqref{eq:5.5.23}, Eq.~\eqref{eq:31.3.3} may be written
\begin{equation}
  \sum_{s=-3/2}^{3/2}
  \langle0|\psi^i(0)|s\rangle
  \langle0|\psi^j(0)|s\rangle^*
  =(2\pi)^{-3}\left(\frac{1+\beta}{2}\right)
  \left[\delta_{ij}-\frac13\gamma_i\gamma_j\right].
  \tag{31.3.4}\label{eq:31.3.4}
\end{equation}
It follows that the momentum-space propagator for a spin \(3/2\) particle
of mass \(m_g\) takes the form
\begin{equation}
  \Delta^{\mu\nu}(q)
  =\frac{P^{\mu\nu}(q)}{q^2+m_g^2-i\epsilon}\,,
  \tag{31.3.5}\label{eq:31.3.5}
\end{equation}
where \(P^{\mu\nu}(q)\) is a Lorentz-covariant polynomial in the four-vector
\(q\), subject to the conditions that for \(\boldsymbol q=0\) and
\(q^0=m_g\), we have
\begin{equation}
  P^{ij}
  =\left(\frac{1+\beta}{2}\right)
   \left[\delta_{ij}-\frac13\gamma_i\gamma_j\right],
  \qquad
  P^{i0}=P^{0i}=P^{00}=0\,.
  \tag{31.3.6}\label{eq:31.3.6}
\end{equation}
Apart from possible terms that vanish on the mass shell (and whose effect
therefore would be the same as direct current--current interactions) the
unique covariant function with this limit is
\begin{equation}
  \begin{aligned}
  P^{\mu\nu}(q)={}&
  \left(\eta^{\mu\nu}+\frac{q^\mu q^\nu}{m_g^2}\right)
  (-i\sl{q}+m_g)
  \\
  &-\frac13
  \left(\gamma^\mu-i\frac{q^\mu}{m_g}\right)
  (i\sl{q}+m_g)
  \left(\gamma^\nu-i\frac{q^\nu}{m_g}\right).
  \end{aligned}
  \tag{31.3.7}\label{eq:31.3.7}
\end{equation}
(The difference between Eq.~\eqref{eq:31.3.7} and any other covariant
function with the limit \eqref{eq:31.3.6} would be a covariant function
whose components all vanish for \(\boldsymbol q=0\) and \(q^0=m_g\) and
hence vanish everywhere on the mass shell.) The free-field Lagrangian
density is then of the form
\begin{equation}
  \mathcal L_0
  =-\frac12\bigl(\bar\psi^\mu D_{\mu\nu}(-i\partial)\psi^\nu\bigr),
  \tag{31.3.8}\label{eq:31.3.8}
\end{equation}
where
\begin{equation}
  \Delta^{\mu\nu}(q)D_{\nu\lambda}(q)=\delta^\mu{}_\lambda\,.
  \tag{31.3.9}\label{eq:31.3.9}
\end{equation}
A tedious but straightforward calculation gives
\begin{equation}
  D_{\nu\lambda}(q)
  =-\epsilon_{\nu\mu\kappa\lambda}\gamma_5\gamma^\mu q^\kappa
   -\frac12m_g[\gamma_\nu,\gamma_\lambda]\,,
  \tag{31.3.10}\label{eq:31.3.10}
\end{equation}
so that the Lagrangian density \eqref{eq:31.3.8} is
\begin{equation}
  \mathcal L_0
  =-\frac12i\epsilon^{\nu\mu\kappa\lambda}
    \bigl(\bar\psi_\nu\gamma_5\gamma_\mu
    \partial_\kappa\psi_\lambda\bigr)
   +\frac14m_g
    \bigl(\bar\psi_\nu[\gamma^\nu,\gamma^\lambda]\psi_\lambda\bigr).
  \tag{31.3.11}\label{eq:31.3.11}
\end{equation}
For \(m_g=0\) this result verifies that the term
\(-\frac12\bar\psi_\mu L^\mu\) in Eq.~\eqref{eq:31.2.9} is the correct
conventionally normalized free-field Lagrangian for a massless
self-charge-conjugate particle of spin \(3/2\). In the limit \(m_g\to0\)
the propagator given by Eqs.~\eqref{eq:31.3.5} and \eqref{eq:31.3.7} is
singular (which just reflects the impossibility shown in Section 5.9 of
forming a \((1,\frac12)+(\frac12,1)\) field from the creation and
annihilation operators for a massless particle of helicity \(\pm3/2\)),
but the terms in Eq.~\eqref{eq:31.3.7} for \(P^{\mu\nu}(q)\) that blow up
for \(m_g\to0\) are all proportional to \(q^\mu\) and/or \(q^\nu\), and
hence give no contribution when the current with which \(\psi_\mu\)
interacts is conserved.

As a further check on the validity of the Lagrangian density
\eqref{eq:31.3.11}, including the peculiar-looking mass term, we note that
it yields a field equation
\begin{equation}
  -i\epsilon^{\nu\mu\kappa\lambda}\gamma_5\gamma_\mu
  \partial_\kappa\psi_\lambda
  +\frac12m_g[\gamma^\nu,\gamma^\lambda]\psi_\lambda=0\,.
  \tag{31.3.12}\label{eq:31.3.12}
\end{equation}
Taking the divergence of this equation yields the result that
\[
  [\sl{\partial},\gamma^\lambda]\psi_\lambda=0\,.
\]
Also, contracting Eq.~\eqref{eq:31.3.12} with \(\gamma_\nu\) shows that
\[
  \gamma_\nu\psi^\nu
  \propto\epsilon^{\nu\mu\kappa\lambda}
  \gamma_5\gamma_\nu\gamma_\mu\partial_\kappa\psi_\lambda
  \propto[\sl{\partial},\gamma^\lambda]\psi_\lambda=0\,,
\]
so that the irreducibility condition \eqref{eq:31.3.1} is satisfied for
free fields (though not necessarily when interactions are taken into
account). From these two results there follows another irreducibility
condition
\[
  \partial_\lambda\psi^\lambda
  =\frac12\{\sl{\partial},\gamma_\lambda\}\psi^\lambda
  =\frac12[\sl{\partial},\gamma_\lambda]\psi^\lambda=0\,.
\]
Using these irreducibility results allows the field equation
\eqref{eq:31.3.12} to be put in the form of a Dirac equation
\begin{equation}
  (\sl{\partial}+m_g)\psi^\lambda=0\,,
  \tag{31.3.13}\label{eq:31.3.13}
\end{equation}
which among other things shows that this is the free field of a particle
of mass \(m_g\).

We will now consider the effects in supergravity theories of a spontaneous
breakdown of supersymmetry. Broken global supersymmetry entails the
existence of a massless particle of spin \(1/2\), the goldstino, but in
supergravity theories the goldstino field \(\chi\) can be eliminated by a
gauge transformation
\(\psi_\mu\to\psi_\mu-\partial_\mu\chi\). Since the gauge is then fixed by
the condition that the goldstino field is eliminated in this way, gauge
invariance no longer keeps the gravitino massless, and it acquires a mass
(to be denoted \(m_g\) from now on), in much the same way as we saw in
Section 21.3 that the vector bosons \(W^\pm\) and \(Z^0\) get masses from
the spontaneous breakdown of the
\(SU(2)\mathbin{\cdot}U(1)\) gauge symmetry of the electroweak interactions.

As discussed at the beginning of Chapter 28, if supersymmetry is relevant
at all to accessible phenomena then the characteristic energy scale at
which it is broken must be much less than the Planck mass. In this case we
may use a continuity argument to give a universal formula for the
gravitino mass \(m_g\). According to Eq.~\eqref{eq:31.1.34}, the gravitino
field \(\psi_\mu\) couples to the supersymmetry current \(S^\mu\) with a
coupling constant \(\frac12\kappa=\frac12\sqrt{8\pi G}\), so the exchange
of a virtual gravitino of four-momentum \(q\) in a transition
\(A+B\to C+D\) contributes to the invariant amplitude a term
\begin{equation}
  \begin{aligned}
  M(A+B\to C+D)
  ={}&\frac14(8\pi G)
  {}_{\mathrm{out}}\!\bra{C}\bar S_\mu\ket{A}_{\mathrm{in}}\!{}_{N}
  \Delta^{\mu\nu}(q)
  {}_{\mathrm{out}}\!\bra{D}S_\nu\ket{B}_{\mathrm{in}}\!{}_{N}\,,
  \end{aligned}
  \tag{31.3.14}\label{eq:31.3.14}
\end{equation}
where the subscript \(N\) indicates that the one-goldstino pole at
\(q^2=0\) has been removed from the matrix element of the supersymmetry
current. For a supersymmetry-breaking scale sufficiently small compared
with the Planck mass, there will be a range of momentum transfers that are
much larger than the gravitino mass but much less than the Planck mass.
For such momenta, the matrix element is dominated by the \(1/m_g^2\)
terms in the propagator numerator \eqref{eq:31.3.7}
\begin{equation}
  \begin{aligned}
  M(A+B\to C+D)\to{}&
  \frac14(8\pi G)
  {}_{\mathrm{out}}\!\bra{C}\bar S_\mu\ket{A}_{\mathrm{in}}\!{}_{N}
  \left(\frac{-2i\sl{q}\,q^\mu q^\nu}{3m_g^2q^2}\right)
  {}_{\mathrm{out}}\!\bra{D}S_\nu\ket{B}_{\mathrm{in}}\!{}_{N}\,.
  \end{aligned}
  \tag{31.3.15}\label{eq:31.3.15}
\end{equation}
But for sufficiently large Planck mass the coupling of the gravitino
becomes negligible, and the matrix element must be the same as would have
been produced by goldstino exchange in a theory without gravitinos.
According to Eq.~\eqref{eq:29.2.10}, this is
\begin{equation}
  \begin{aligned}
  M(A+B\to C+D)\to{}&
  {}_{\mathrm{out}}\!\bra{C}\bar S_\mu\ket{A}_{\mathrm{in}}\!{}_{N}
  \left(\frac{-i\sl{q}}{q^2}\right)
  \left(\frac{q^\mu q^\nu}{F^2}\right)
  {}_{\mathrm{out}}\!\bra{D}S_\nu\ket{B}_{\mathrm{in}}\!{}_{N}\,,
  \end{aligned}
  \tag{31.3.16}\label{eq:31.3.16}
\end{equation}
where \(F\) is the parameter characterizing the strength of supersymmetry
breaking (here taken real), defined so that the vacuum energy density is
\(F^2/2\). In order for Eqs.~\eqref{eq:31.3.15} and
\eqref{eq:31.3.16} to agree, the gravitino mass must have the value
\begin{equation}
  m_g=\sqrt{\frac{4\pi G F^2}{3}}\,.
  \tag{31.3.17}\label{eq:31.3.17}
\end{equation}
This formula is valid only to lowest order in \(GF^2\), but to all orders
(and even non-perturbatively) in the non-gravitational interactions
responsible for the spontaneous breakdown of supersymmetry.

It is convenient for some purposes to express \(m_g\) in terms of the
expectation values \(\langle s\rangle\) and \(\langle p\rangle\) of the
spinless auxiliary gravitational fields. Note that for the vacuum state
to have zero spacetime curvature, the vacuum energy density \(F^2/2\) of
matter fields must be balanced by the negative vacuum energy of
gravitation and its interaction with the hidden sector fields, which is
given in terms of \(\langle s\rangle\) and \(\langle p\rangle\) by
Eqs.~\eqref{eq:31.2.18} and \eqref{eq:31.2.16} as
\(-\frac43(\langle s\rangle^2+\langle p\rangle^2)\), so
\begin{equation}
  F^2/2=\frac43\bigl(\langle s\rangle^2+\langle p\rangle^2\bigr).
  \tag{31.3.18}\label{eq:31.3.18}
\end{equation}
We can therefore write Eq.~\eqref{eq:31.3.17} as
\begin{equation}
  m_g=\frac{2\kappa}{3}
  \sqrt{\langle s\rangle^2+\langle p\rangle^2}\,.
  \tag{31.3.19}\label{eq:31.3.19}
\end{equation}
It is sometimes convenient to introduce a complex gravitino mass, defined
as
\begin{equation}
  \widetilde m_g\equiv
  \frac{2\kappa}{3}\bigl(\langle s\rangle+i\langle p\rangle\bigr),
  \tag{31.3.20}\label{eq:31.3.20}
\end{equation}
whose absolute magnitude is the physical gravitino mass
\eqref{eq:31.3.19}.
